% SPDX-License-Identifier: GPL-3.0-or-later
% Copyright (C) 2026 Alessandro Greco
\chapter{运行中继或网桥}\label{cap:22-relay}
\aggiornatoal{2026 年 9 月 23 日}{tor 0.4.9.13（手册）；实验网络}

\begin{obiettivi}
学完本章后，你将能够：
\begin{itemize}[nosep]
  \item 根据自己的资源和能够承担的风险，选择运行中间中继、出口中继或网桥；
  \item 编写各类节点所需的基本配置，包括带宽限制与离线密钥存储；
  \item 阅读、编写出口策略，并在实验网络中验证其作用（实验~14）；
  \item 正确声明中继家族。
\end{itemize}
\end{obiettivi}

Tor 网络之所以存在，是因为数千名个人和组织贡献了带宽与计算机。前面的章节从客户端视角观察中继；现在我们换到运营者的一侧。本章并不是系统管理手册——Tor Project 已维护一份持续更新的运营指南——而是解释主要的运维选择在协议层面意味着什么。

% ---------------------------------------------------------------------------
\section{应该运行哪种节点？}\label{sec:relay-tipi}

\begin{description}
  \item[中间中继]（之后也可能成为入口守卫）在其他中继之间转发 cell，但不向互联网外部目的地建立连接。这是运营要求最低的一种角色。其地址会出现在共识中，但网站不会看到来自它的出口流量。
  \item[出口中继] 代表客户端向目标网站或服务建立连接。由于出口相对于需求仍然稀缺（\cref{sec:percorso-pesi}），这一角色尤其有价值，但也带来最高的风险暴露：如果出现滥用流量，接收方看到的表面来源就是出口的 IP 地址。
  \item[网桥] 不列入公开列表的中继，供无法直接访问 Tor 公共网络的用户作为第一跳（\cref{sec:cens-bridge}）。即使只是普通住宅宽带，只要配合 obfs4 或 WebTunnel 等可插拔传输，也能有所帮助。
\end{description}

\begin{attenzione}[开始之前]
运行出口中继会产生法律和运营后果，包括滥用投诉以及有关部门的查询，参见 \cref{cap:25-legale}。可以考虑先从中间中继或网桥开始，事先告知网络接入服务商，并避免在同一条网络连接上既运行中继，又以 Tor 客户端身份开展敏感活动。
\end{attenzione}

% ---------------------------------------------------------------------------
\section{基本配置}\label{sec:relay-config}

中间中继只需要在 \file{torrc} 中设置几个选项：
\begin{lstlisting}[numbers=none]
Nickname        MyRelay
ContactInfo     operator@example.org
ORPort          443
ExitRelay       0
SocksPort       0
RelayBandwidthRate  4 MBytes
RelayBandwidthBurst 8 MBytes
\end{lstlisting}
\opt{ContactInfo} 用于在中继出现问题时让其他人联系运营者。手册提醒，描述符会被归档和公开，因此联系地址可能无限期地保留在公共记录中。\opt{ExitRelay 0} 明确禁止该节点作为出口；对于近年的版本，这本已是默认行为：设置为 \texttt{auto} 时，只有运营者编写出口策略，tor 才会成为出口。\opt{RelayBandwidthRate} 和 \opt{RelayBandwidthBurst} 分别限制平均与峰值带宽。相比之下，\opt{AccountingMax} 限制一个统计周期（例如一个月）内的总流量。一旦达到上限，中继便停止接受新流量，直到下一个周期开始。

\textbf{运维注意：}当 \texttt{ExitRelay} 设为 \texttt{auto} 时，不仅显式的 \texttt{ExitPolicy}，\texttt{ReducedExitPolicy} 或 \texttt{IPv6Exit} 配置也可能使节点成为出口。示例中的 \texttt{ExitRelay 0} 则会禁止全部出口流量。部署公开中继之前应检查最终生效的配置。

\subsection{新中继的成长过程}

刚启动的中继会向目录权威节点发布描述符，随后由权威节点将其纳入共识并开始测量（\cref{sec:dir-voto}）。起初，它获得的流量很少：客户端根据带宽权重选择中继，测量结果稳定前，新节点的权重仍然较低。各种标志需要随着运行时间逐渐获得（\cref{sec:dir-flag}）。\texttt{Stable} 需要长期持续、不中断地运行，而 \texttt{Guard} 还要求节点达到一定年龄并保持高可用性。当一台中继刚成为入口守卫时，它的流量甚至可能暂时下降：已经选好守卫的客户端，不会像过去那样频繁把它用作中间中继。随着新客户端逐渐选择它，流量才会缓慢增长（\cref{sec:guard}）。

两种故障的后果并不相同。频繁重启但不改变身份的中继，仍保留带宽测量数据，却会降低获得 \texttt{Stable} 和 \texttt{Guard} 所需的持续可用性，因此可能需要重新取得这些标志。相反，改变\emph{身份密钥}会让目录权威节点将其视为一个全新的中继，就连带宽测量也必须从零开始。

\subsection{密钥}

中继的身份由其主 Ed25519 密钥决定（\cref{sec:canale-chiavi}）。默认情况下，tor 将它存放在中继磁盘上。如果设置 \opt{OfflineMasterKey 1}，就可以把主密钥保留在另一台计算机上，只把使用 \cmd{tor --keygen} 签发证书的临时签名密钥传输给服务器。即使服务器遭到攻破，攻击者获得的也是会过期的密钥，而不是中继的长期身份。不过这种方法需要严格管理：临时密钥必须在到期前更新，否则中继将停止正常运行。

% ---------------------------------------------------------------------------
\section{出口策略}\label{sec:relay-exitpolicy}

出口策略是一份按顺序排列的 \texttt{accept} 与 \texttt{reject} 规则列表。每条规则指定一个地址或网络，以及一个端口或端口范围。Tor \textbf{按顺序}检查规则，并采用第一条匹配规则。运营者自行定义的规则优先于手册列出的默认规则。默认情况下，Tor 会拒绝用于电子邮件的端口（25）、NNTP（119）、NetBIOS 与 SMB（135--139、445），以及部分文件共享应用使用的端口，其余流量则允许通过。另外，tor 默认拒绝私有目的地址，例如 \texttt{127.0.0.0/8}、\texttt{10.0.0.0/8}、\texttt{192.168.0.0/16}，以及中继自身的公开地址。这可以保护运营者的内部网络（\opt{ExitPolicyRejectPrivate}）。我们的实验网络关闭了这一保护，以允许流量连接到 \texttt{127.0.0.1}。

\textbf{保护规则的先后顺序：}默认开启的 \texttt{ExitPolicyRejectPrivate} 会把拒绝私有网络和本地地址的规则\emph{放在}运营者自定义规则之前。Chutney 私有实验网络故意关闭了这项保护；不要将这种设置照搬到公开出口中继。

只允许 Web 流量通过的出口可以配置为：
\begin{lstlisting}[numbers=none]
ExitRelay  1
ExitPolicy accept *:80
ExitPolicy accept *:443
ExitPolicy reject *:*
\end{lstlisting}
Tor Project 还建议将 \opt{ReducedExitPolicy} 作为一种常见折中方案：它允许若干广泛使用的服务端口，但排除更容易产生滥用投诉的端口。

\begin{figure}[htbp]
\centering
% SPDX-License-Identifier: GPL-3.0-or-later
% Copyright (C) 2026 Alessandro Greco
%% EPUB 无障碍替代文本。
\altfig{Chutney 实验出口策略：只在测试网络中设置 ExitPolicyRejectPrivate 0，关闭对私有地址的优先拒绝。两个发送到 127.0.0.1 的端口 8080 与 8081 的 BEGIN 请求按顺序经过规则。8080 命中 reject 127.0.0.1:8080，返回 END EXITPOLICY；8081 不命中拒绝，而命中 accept 127.0.0.0/8:*，因此可以尝试 TCP 连接，但只有服务器监听时才可能收到 CONNECTED。若把 accept 提前，特定拒绝就无效。公开出口不得禁用 ExitPolicyRejectPrivate。}
% 仅在实验室中的出口策略
\begin{tikzpicture}[font=\sffamily\scriptsize,node distance=0.45cm,
  reg/.style={draw,rounded corners=1pt,minimum height=0.6cm,text width=4.6cm,font=\ttfamily\scriptsize},
  esito/.style={draw,rounded corners=2pt,minimum height=0.6cm,text width=2.5cm,align=center}]
  \node[font=\sffamily\scriptsize] (top) at (0,1.25) {两个 BEGIN：127.0.0.1:8080 / :8081};
  \node[draw,rounded corners=2pt,fill=torviola!10,font=\ttfamily\tiny] (r0) at (0,0.15) {ExitPolicyRejectPrivate 0};
  \node[reg,fill=avviso!12,below=of r0] (r1) {reject 127.0.0.1:8080};
  \node[reg,fill=torverde!12,below=of r1] (r2) {accept 127.0.0.0/8:*};
  \node[reg,fill=torgrigio!10,below=of r2] (r3) {（其他默认策略）};
  \draw[->] (top) -- (r0); \draw[->] (r0) -- (r1); \draw[->] (r1) -- node[right,font=\sffamily\tiny] {8081：不匹配} (r2); \draw[->] (r2) -- node[right,font=\sffamily\tiny] {其他情况} (r3);
  \node[esito,fill=avviso!15,draw=avviso,right=1.2cm of r1] (e1) {\texttt{END} 原因\\\texttt{EXITPOLICY}};
  \draw[->,avviso,thick] (r1) -- node[above,font=\sffamily\tiny] {8080 匹配} (e1);
  \node[esito,fill=torverde!15,draw=torverde!60!black,right=1.2cm of r2] (e2) {尝试 TCP\\服务器监听时\\才可 CONNECTED};
  \draw[->,torverde!60!black,dashed] (r2) -- node[above,font=\sffamily\tiny] {8081 匹配} (e2);
  \node[anchor=west,font=\sffamily\tiny,text=torgrigio,align=left] at (-5.9,-0.55) {第一条匹配规则\\决定结果};
  \node[anchor=west,font=\sffamily\tiny,text=avviso] at (-5.9,-1.95) {不可用于公开出口};
  \node[font=\sffamily\tiny,text=torgrigio,anchor=south] at (r0.north) {仅实验：ExitPolicyRejectPrivate 0};
\end{tikzpicture}

\caption{Chutney 实验室仅为连接私有地址而设置 \opt{ExitPolicyRejectPrivate 0}；正常情况下 Tor 会在运营者自定义规则之前拒绝这些地址。到 \texttt{127.0.0.1:8080} 的 \texttt{BEGIN} 命中 \texttt{reject} 并产生 \texttt{END EXITPOLICY}；端口 8081 跳过该规则并命中 \texttt{accept}。允许仅代表可以尝试 TCP 连接，服务器实际监听时才可能收到 \texttt{CONNECTED}。颠倒规则顺序会使拒绝无效。}
\label{fig:policy}
\end{figure}

\paragraph{客户端得到什么信息？}中继描述符中包含完整的出口策略，而共识和微描述符中只提供按端口汇总的摘要，例如 \texttt{p accept 80,443}。客户端借此判断某个出口能否处理请求（\cref{sec:percorso-vincoli}）。因为该摘要省略了目的地址，客户端可能只有在尝试建立连接后，才能发现某个地址受到特定规则的拒绝，正如下面的实验所示。

% ---------------------------------------------------------------------------
\section{实验：测试出口策略}\label{sec:lab-relay}

\begin{laboratorio}[14]{出口规则及其后果}
\textbf{目标。}给实验出口增加一条规则，再用 \file{labs/zh-Hans/代码/torlab} 观察允许和拒绝不同目的地时的回应。\\
\textbf{前置条件。}实验~3 和~4；实验网络已经运行。\\
\textbf{预计时间。}1 小时。
\end{laboratorio}

\paragraph{步骤 1：添加规则。}在 \texttt{test006e} 的 \file{torrc} 中，在原先允许 \texttt{127.0.0.0/8} 的规则前面加入一条拒绝某个端口的规则：
\begin{lstlisting}[numbers=none]
ExitPolicy reject 127.0.0.1:8080
ExitPolicy accept 127.0.0.0/8:*
\end{lstlisting}
这里的两行配置是相对顺序示意：修改已有的 \file{torrc} 时，只需把新的 \texttt{reject} 插入第一条 \texttt{ExitPolicy} 之前，切勿重复添加原有的 \texttt{accept} 规则。顺序非常重要。如果把拒绝规则放到 \texttt{accept} 后面，它就不会产生作用，因为第一条匹配规则已经决定了结果。在之前的参考实验中，我们曾犯过这个错误。重新启动出口中继后，在容器的 8080 和 8081 端口各启动一个简单的 HTTP 服务器。

\paragraph{步骤 2：发送三个请求。}使用 \texttt{build} 和 \file{stream.py} 建立一条以 \texttt{006e} 为出口的电路，分别向端口 8080、8081，以及没有服务监听的 9999 发送 \texttt{BEGIN} 消息。在我们的一次会话中，响应如下（原因码、地址和 TTL 可能随版本与环境变化）：
\begin{lstlisting}[numbers=none]
127.0.0.1:8080 -> END  reason=EXITPOLICY address=127.0.0.1 TTL=1800
127.0.0.1:8081 -> CONNECTED
127.0.0.1:9999 -> END  reason=MISC
\end{lstlisting}

\paragraph{步骤 3：解释回应。}策略拒绝会产生 \texttt{EXITPOLICY} 原因码。\texttt{END} cell 还会包含解析后的地址和 TTL。客户端因此可以记住该出口不能访问这个目的地，转而选择另一个出口，而无须等待新描述符。相比之下，没有服务器监听的端口返回通用的 \texttt{MISC}，而不是规范针对连接被拒绝所列出的 \texttt{CONNECTREFUSED}~\spec{tor-spec}{Closing streams}。

\begin{approfondimento}[为什么得到 \texttt{MISC} 而不是 \texttt{CONNECTREFUSED}？]
答案出现在出口中继的 \texttt{info} 日志中。追踪这一问题，可以看到研究规范与实际实现存在差异时应如何调查：
\begin{lstlisting}[numbers=none]
[info] connection_handle_write_impl(): in-progress connect failed.
       Removing. (Connection refused)
[info] errno_to_stream_end_reason(): Didn't recognize errno 115 (Operation
       now in progress); telling the client that we are ending a stream
       for 'misc' reason.
\end{lstlisting}
出口中继实际上\emph{知道}连接被拒绝，却依据另一个数值确定原因码。到目的地的连接采用非阻塞方式：\texttt{connect()} 最初返回 \texttt{EINPROGRESS}（errno 115），实际结果要等后续写事件发生时才会得知。此时，\file{connection.c} 通过 \texttt{getsockopt(SO\_ERROR)} 获取真正的错误，记录 \enquote{Connection refused}，随后调用 \file{connection\_edge\_end\_errno()}。然而，该函数又使用 \texttt{errno\_to\_stream\_end\_reason(tor\_socket\_errno(\ldots))} 重新计算原因码。在 Unix 上，\texttt{tor\_socket\_errno} 是展开为 \texttt{(errno)} 的宏（\file{src/lib/net/socket.h}），读取的仍是保留了 \texttt{EINPROGRESS} 的全局变量。数值 115 不在 \file{src/core/or/reasons.c} 的已识别列表中，因此默认分支返回 \texttt{MISC}。

正确信息本已存在，却没有被使用。相关代码中的两个分支也有所区别：中继到中继的连接分支会传递刚获取的错误，而流处理分支则忽略它。根据已有的 \file{docs/it/segnalazioni/} 分析与补丁记录，此前在私有网络中对 tor 0.4.9.13 进行的测试表明，正确传递错误值后第三个请求返回 \texttt{CONNECTREFUSED}。这是一个不影响匿名性的轻微缺陷：客户端仍能知道流建立失败，但失去了更有用的信息。\texttt{CONNECTREFUSED} 会提示目的地明确拒绝了连接，因此换一个出口重试也可能没有意义。分析报告与补丁存放在 \file{docs/it/segnalazioni/}，可供 Tor Project 审阅。这也说明为什么必须用真实行为检验协议规范。
\end{approfondimento}

实际脚本为 \file{labs/zh-Hans/实验室/exit\_policy.py}。在 Linux 宿主机的专用 Chutney 网络启动后，可从仓库根目录执行 \texttt{PYTHONPATH=labs/zh-Hans/代码 python3 labs/zh-Hans/实验室/exit\_policy.py}。该脚本运行三项测试，并在 \texttt{finally} 中尝试恢复 \texttt{006e} 的原始 \file{torrc}。它会暂时修改出口配置并开启本地 HTTP 服务器，\emph{只能用于可丢弃的私有实验网络}。常规的 \texttt{make -C labs/zh-Hans/实验室 verify} 只执行 Chutney 检查，并不会调用此脚本。本轮审核未运行脚本，也未动态验证恢复过程。有一个重要细节：电路必须包含三跳。出口会拒绝转发来自单跳电路的流量，并返回原因码 \texttt{TORPROTOCOL} 的 \texttt{END}（\cref{sec:dos-limiti}）。

% ---------------------------------------------------------------------------
\section{中继家族}\label{sec:relay-famiglie}

运营多个中继的人必须声明一个\textbf{家族}（family）。客户端不会将同一家族中的两个中继放入同一条电路（\cref{sec:percorso-vincoli}），因为同一运营者可能关联所观察到的流量。传统机制是 \opt{MyFamily}：每台中继都列出其他所有中继的指纹，只有双方声明互相认可时，家族关系才成立。如果维护十台中继，就需要管理十份一致的列表，每次加入新中继时都要修改它们。

tor 0.4.9 引入的新机制使用\textbf{家族密钥}。运营者通过 \cmd{tor --keygen-family MyFamilyKey} 生成密钥，将生成的 \texttt{MyFamilyKey.secret\_family\_key} 文件复制到每台中继的密钥目录（\texttt{KeyDirectory}，文件名须保留 \texttt{.secret\_family\_key} 后缀），并在每台中继的 \file{torrc} 中写入生成的 \texttt{FamilyId} 配置行。必须保护并备份这份共享的私钥；若密钥泄露，应为整个家族轮换密钥。每台中继都用该密钥签署把自身身份绑定到家族的证书。验证方因而无须逐一比较列表。在过渡阶段，建议同时配置传统机制和新机制。

\begin{approfondimento}[为什么声明家族有帮助？]
没有什么强制运营者把自己的中继声明为同一家族；攻击者甚至可能故意不这样做。不过，监测网络的目录权威节点和志愿者会留意共同特征，例如相同的网络、软件版本、同时启动的时间和联系信息。未声明的家族可能引起额外审查，相关中继甚至可能被排除在共识之外。主动声明家族，是整个网络信任安排的一部分。
\end{approfondimento}

% ---------------------------------------------------------------------------
\textbf{不要将网桥列入公开的 \texttt{MyFamily} 列表：}这可能暴露网桥与公开中继的关联。Tor 官方运维指南明确建议不要在网桥上设置该选项。此外，\texttt{tor --keygen-family MyFamilyKey} 可能覆盖同名密钥文件；请使用未被占用的名称，并严格保护生成的秘密文件。

\section{运行网桥}\label{sec:relay-bridge}

可以按照实验~10（\cref{sec:lab-censura}）配置一台 obfs4 网桥：
\begin{lstlisting}[numbers=none]
BridgeRelay 1
ORPort 9001
ServerTransportPlugin obfs4 exec /usr/bin/obfs4proxy
ServerTransportListenAddr obfs4 0.0.0.0:443
ExtORPort auto
BridgeDistribution any
\end{lstlisting}
\opt{BridgeDistribution} 告诉当前的分发系统 rdsys，运营者希望如何分发网桥。\texttt{any} 将选择交给系统，\texttt{none} 则完全排除公开分发，适合仅向可信的人私下分享网桥。网桥信息页面列出了可用渠道，包括网站、电子邮件和 Telegram。网桥虽然不出现在公开共识中，仍会把描述符发送给网桥权威节点，由后者检查可达性。

传输程序应监听不同于 ORPort 的端口；条件允许时，最好使用 443 等常见端口。这才是客户端使用、也是审查者观察到的端口（\cref{sec:cens-strategie}）。

% ---------------------------------------------------------------------------
\section*{本章小结}
\begin{itemize}
  \item 运行中间中继是门槛最低的贡献方式；出口中继特别有价值，但风险暴露更高；网桥可以帮助面临审查的用户。
  \item 基本配置只需要身份、联系信息、端口和带宽限制等少数选项。主身份密钥可以离线保存。
  \item 出口策略是按顺序排列的规则列表，第一条匹配规则决定结果；默认拒绝私有目的地址。
  \item 策略拒绝以 \texttt{END} 消息传给客户端，包含 \texttt{EXITPOLICY} 原因码、地址和 TTL。
  \item 多中继运营者应使用 \opt{MyFamily} 或较新的家族密钥机制声明家族关系。
\end{itemize}

% ---------------------------------------------------------------------------
\section*{练习}

\begin{esercizio}[1]\label{es:22-ordine}
给定策略 \texttt{accept *:80}、\texttt{reject 192.0.2.0/24:*}、\texttt{reject *:*}，客户端能否通过该出口连接 \texttt{192.0.2.7:80}？请修改策略，使这条连接被拒绝。
\end{esercizio}

\begin{esercizio}[1]\label{es:22-riassunto}
为什么共识中仅包含端口信息的摘要，不足以判断出口是否允许某次连接？客户端收到原因码为 \texttt{EXITPOLICY} 的 \texttt{END} 后会怎么做？
\end{esercizio}

\begin{esercizio}[2]\label{es:22-famiglia}
某运营者使用 \opt{MyFamily} 管理五台中继，又加入第六台，却只更新了新中继的配置。客户端会看到什么？改用家族密钥后有什么变化？
\end{esercizio}

\begin{esercizio}[2]\label{es:22-refused}
验证技术分析中的解释：在出口中继启用 \texttt{info} 日志，重新执行三个请求。如果 \file{connection\_edge\_end\_errno()} 收到的是 \texttt{getsockopt} 返回的错误值，会发生什么变化？哪些其他流关闭原因也可能因此被正确报告？
\end{esercizio}

\begin{esercizio}[3]\label{es:22-accounting}
一台中继的月流量上限为 \texttt{AccountingMax 1 TBytes}，带宽为 10\,MB/s。全速运行多少天后会耗尽配额？一台在半个月内处于离线状态的中继，与另一台整月在线但带宽减半的中继相比，会对网络产生什么不同影响？
\end{esercizio}

% ---------------------------------------------------------------------------
\section*{延伸阅读}
\begin{itemize}
  \item Tor Project 的中继运营指南：\url{https://community.torproject.org/relay/}。
  \item \begin{sloppypar}tor 手册中的 \opt{ExitPolicy}、\opt{ReducedExitPolicy}、
        \opt{FamilyId}、\opt{OfflineMasterKey} 与 \opt{BridgeDistribution} 条目。\end{sloppypar}
  \item 流关闭原因规范：\spec{tor-spec}{Closing streams}。
\end{itemize}
